\documentclass[11pt,a4paper]{article}
\usepackage[margin=1in]{geometry}
\usepackage{amsmath,amssymb,amsthm}
\usepackage{algorithm}
\usepackage{algpseudocode}
\usepackage{booktabs}
\usepackage{hyperref}

\theoremstyle{definition}
\newtheorem{definition}{\textbf{Definition}}[section]
\newtheorem{theorem}{\textbf{Theorem}}[section]
\newtheorem{lemma}[theorem]{\textbf{Lemma}}
\newtheorem{remark}{\textbf{Remark}}[section]

\title{\textbf{Soft Actor-Critic (SAC)}}
\author{\textbf{Team Scaibu} \\ \small Email: \texttt{jaydeep.9664920749@gmail.com} \\ \small \textbf{Architecture and Core Engine Team}}
\date{\textbf{October 2026}}

\begin{document}
\maketitle

\section{Definitions and Cognitive Mental Models}

\subsection{Formal Mathematical Definition}
\textbf{Definition (Maximum Entropy Markov Decision Processes and Soft Bellman Operators):}
Standard reinforcement learning maximizes the expected cumulative reward $J(\pi) = \sum_{t} \mathbb{E}_{(s_t, a_t) \sim \rho_\pi} [\gamma^t r(s_t, a_t)]$. In complex continuous action spaces, this objective frequently leads to premature convergence to brittle, suboptimal deterministic policies. \textbf{Soft Actor-Critic (SAC)} generalizes the standard formulation to the \textbf{Maximum Entropy Reinforcement Learning} framework by augmenting the objective with an explicit policy entropy regularization term:
{\boldmath
\begin{equation}
J_{\text{MaxEnt}}(\pi) = \sum_{t=0}^{\infty} \mathbb{E}_{(s_t, a_t) \sim \rho_\pi} \left[ \gamma^t \left( r(s_t, a_t) + \alpha \mathcal{H}(\pi(\cdot | s_t)) \right) \right]
\end{equation}
}
where $\mathcal{H}(\pi(\cdot | s_t)) = -\mathbb{E}_{a_t \sim \pi} \left[ \ln \pi(a_t | s_t) \right]$ represents the differential entropy of the action distribution, and $\alpha > 0$ is the stochastic temperature parameter determining the relative trade-off between reward exploitation and policy entropy exploration.

To evaluate state-action values under the soft objective, SAC employs the \textbf{Soft Bellman Expectation Operator} $\mathcal{T}^\pi$:
{\boldmath
\begin{equation}
\mathcal{T}^\pi Q(s_t, a_t) = r(s_t, a_t) + \gamma \mathbb{E}_{s_{t+1} \sim p} \left[ V(s_{t+1}) \right]
\end{equation}
}
where the soft state-value function $V(s_{t+1})$ satisfies the analytical identity:
{\boldmath
\begin{equation}
V(s_{t+1}) = \mathbb{E}_{a' \sim \pi} \left[ Q(s_{t+1}, a') - \alpha \ln \pi(a' | s_{t+1}) \right]
\end{equation}
}
To prevent continuous positive maximization bias across policy iterations, SAC maintains clipped double Q-networks $Q_{\theta_1}$ and $Q_{\theta_2}$, yielding the target:
{\boldmath
\begin{equation}
y_t = r(s_t, a_t) + \gamma \left( \min_{i \in \{1, 2\}} Q_{\bar{\theta}_i}(s_{t+1}, a') - \alpha \ln \pi_\phi(a' | s_{t+1}) \right), \quad a' \sim \pi_\phi(\cdot | s_{t+1})
\end{equation}
}

\subsection{Expanded Non-Technical Definition (Intuition for Anyone)}
\textbf{Intuitive Conceptual Definition:}
\begin{enumerate}
    \item \textbf{High Reward While Staying as Open-Minded as Possible}: Standard robots try to find one exact path that gives points, quickly ignoring other paths. If a sudden obstacle appears on that path, the robot fails. Soft Actor-Critic instructs the agent: ``Succeed at the task, but maintain the maximum possible randomness and flexibility in how you achieve it!''
    \item \textbf{Built-in Exploration}: Because high entropy is directly rewarded, the agent explores naturally without needing artificial noise injection. It discovers multiple ways to solve problems and avoids getting stuck in local traps.
    \item \textbf{Clipped Double Q-Learning}: The agent uses two independent critics and always trusts the more pessimistic estimate. This stops the policy from over-optimistically chasing random statistical noise in value predictions.
    \item \textbf{Self-Tuning Temperature}: As the agent masters the environment, it automatically adjusts its randomness parameter $\alpha$, transitioning smoothly from wild exploratory discovery to precise expert execution.
\end{enumerate}

\subsection{Child-Friendly Mental Model (Physical Metaphor)}
Imagine an acrobat learning to cross a rope bridge in a windstorm:
\begin{itemize}
    \item \textbf{The Rigid Robot (Traditional RL)}: The acrobat memorizes one single rigid foot placement. If a gust of wind blows, they have no balance adjustments and fall.
    \item \textbf{The Flexible Cat (Soft Actor-Critic)}: The acrobat learns to cross the bridge while keeping all their muscles relaxed, flexible, and ready to bounce or pivot in any direction.
    \item \textbf{Two Honest Spotters}: Two spotters watch each step. Whenever there is disagreement about how stable a step is, the acrobat listens to the more cautious spotter.
    \item \textbf{Effortless Mastery}: Because the acrobat knows 50 different ways to recover their balance, they dance across the rope bridge safely under any storm!
\end{itemize}

\section{Core Mathematical Equations and Definitions}

\subsection{Soft Bellman Value Equation}
{\boldmath
\begin{equation}
V(s_t) = \mathbb{E}_{a_t \sim \pi} \left[ \min_{i \in \{1, 2\}} Q_i(s_t, a_t) - \alpha \ln \pi(a_t | s_t) \right]
\end{equation}
}
\begin{itemize}
    \item \textbf{Formal Definition}: The soft state-value function combining minimum expected Q-value with entropy reward.
    \item \textbf{What It Means}: Evaluates how desirable a state is based on expected future rewards plus the variety of good moves available.
    \item \textbf{Why It Is Important}: Connects policy entropy to dynamic programming value iteration without requiring explicit value network approximation.
\end{itemize}

\subsection{Clipped Soft Bellman Critic Target}
{\boldmath
\begin{equation}
y_t = r_t + \gamma (1 - d_t) \left( \min(Q_{\bar{\theta}_1}(s_{t+1}, a'), Q_{\bar{\theta}_2}(s_{t+1}, a')) - \alpha \ln \pi_\phi(a' | s_{t+1}) \right)
\end{equation}
}
\begin{itemize}
    \item \textbf{Formal Definition}: The temporal difference target combining real transition reward with clipped double target Q-estimates.
    \item \textbf{What It Means}: Uses the lower of two future value predictions minus the entropy cost to train both critic networks via MSE.
    \item \textbf{Why It Is Important}: Prevents positive value overestimation bias that causes standard actor-critic policies to diverge.
\end{itemize}

\subsection{Reparameterized Policy Loss}
{\boldmath
\begin{equation}
J_\pi(\phi) = \mathbb{E}_{s_t \sim \mathcal{D}, \epsilon_t \sim \mathcal{N}(0, I)} \left[ \alpha \ln \pi_\phi(f_\phi(\epsilon_t; s_t) | s_t) - \min_{i \in \{1, 2\}} Q_{\theta_i}(s_t, f_\phi(\epsilon_t; s_t)) \right]
\end{equation}
}
\begin{itemize}
    \item \textbf{Formal Definition}: The policy gradient objective using the pathwise derivative estimator via Gaussian squashed reparameterization $a = \tanh(\mu_\phi(s) + \sigma_\phi(s) \odot \epsilon)$.
    \item \textbf{What It Means}: Backpropagates critic gradients directly through the actor network parameters to maximize soft value.
    \item \textbf{Why It Is Important}: Produces substantially lower gradient variance than score-function REINFORCE estimators, drastically accelerating convergence.
\end{itemize}

\subsection{Dual Entropy Temperature Optimization}
{\boldmath
\begin{equation}
J(\alpha) = \mathbb{E}_{a_t \sim \pi_t} \left[ -\alpha \ln \pi_t(a_t | s_t) - \alpha \bar{\mathcal{H}} \right]
\end{equation}
}
\begin{itemize}
    \item \textbf{Formal Definition}: The convex dual objective function for the entropy temperature parameter subject to minimum entropy constraint $\bar{\mathcal{H}}$.
    \item \textbf{What It Means}: Increases $\alpha$ when policy entropy drops below target $\bar{\mathcal{H}}$, and decreases $\alpha$ when entropy is comfortably high.
    \item \textbf{Why It Is Important}: Eliminates manual tuning of the entropy coefficient across different tasks and continuous control dimensions.
\end{itemize}

\section{Rigorous Mathematical Analysis Checklist}

\begin{itemize}
    \item \textbf{Notation}: $s \in \mathcal{S}$ (continuous state), $a \in \mathcal{A}$ (continuous action), $\alpha$ (temperature), $\gamma \in (0, 1]$ (discount factor), $\bar{\mathcal{H}}$ (target entropy), $Q_1, Q_2$ (twin critics).
    \item \textbf{Epistemic Status}: Formulated by Haarnoja et al. (UC Berkeley, 2018), widely acknowledged as the state-of-the-art continuous control benchmark algorithm.
    \item \textbf{Dependency Graph}: Experience buffer $\to (s, a, r, s', d) \to a' \sim \pi \to y_t \to \nabla_{\theta_1, \theta_2} J_Q \to \text{Reparameterized } a \to \nabla_\phi J_\pi \to \nabla_\alpha J(\alpha)$.
    \item \textbf{Dimensions}: Transition batch size $N$, continuous action dimension $d_A$, state dimension $d_S$.
    \item \textbf{Regimes}: Continuous robotic manipulation, quadrupedal locomotion, autonomous vehicle flight control, industrial process setpoint optimization.
    \item \textbf{Invariants}: If terminal flag $d_t = 1$, soft target $y_t \equiv r_t$ identically, independent of future values or entropy.
    \item \textbf{Isomorphisms}: SAC is isomorphic to continuous soft Q-learning and variational inference under energy-based policy priors.
\end{itemize}

\section{Theoretical Theorems and Formal Proofs}

\begin{theorem}[Soft Policy Improvement Theorem]
Let $\pi_{\text{new}}$ be defined as the Boltzmann distribution proportional to $\exp\left(\frac{1}{\alpha} Q^{\pi_{\text{old}}}(s, \cdot)\right)$. Then for all $(s, a) \in \mathcal{S} \times \mathcal{A}$:
{\boldmath
\begin{equation}
Q^{\pi_{\text{new}}}(s, a) \ge Q^{\pi_{\text{old}}}(s, a)
\end{equation}
}
and the policy improvement inequality holds strictly unless $\pi_{\text{old}}$ is already optimal.
\end{theorem}

\begin{proof}
Let $\mathcal{T}^{\pi}$ be the soft Bellman operator. We have:
{\boldmath
\begin{equation}
Q^{\pi_{\text{old}}}(s, a) = r(s, a) + \gamma \mathbb{E}_{s' \sim p} \left[ V^{\pi_{\text{old}}}(s') \right]
\end{equation}
}
By definition of the soft value function under $\pi_{\text{old}}$:
{\boldmath
\begin{equation}
V^{\pi_{\text{old}}}(s) = \mathbb{E}_{a \sim \pi_{\text{old}}} \left[ Q^{\pi_{\text{old}}}(s, a) - \alpha \ln \pi_{\text{old}}(a | s) \right]
\end{equation}
}
Now consider the Kullback-Leibler divergence between $\pi_{\text{new}}(\cdot | s)$ and the Boltzmann policy:
{\boldmath
\begin{align}
D_{\text{KL}}\left(\pi_{\text{new}}(\cdot | s) \,||\, \frac{\exp\left(\frac{1}{\alpha} Q^{\pi_{\text{old}}}(s, \cdot)\right)}{Z(s)}\right) &\ge 0 \\
\mathbb{E}_{a \sim \pi_{\text{new}}}\left[ \ln \pi_{\text{new}}(a | s) - \frac{1}{\alpha} Q^{\pi_{\text{old}}}(s, a) + \ln Z(s) \right] &\ge 0
\end{align}
}
Multiplying by $-\alpha$ reverses the inequality:
{\boldmath
\begin{equation}
\mathbb{E}_{a \sim \pi_{\text{new}}}\left[ Q^{\pi_{\text{old}}}(s, a) - \alpha \ln \pi_{\text{new}}(a | s) \right] \ge \alpha \ln Z(s) \ge V^{\pi_{\text{old}}}(s)
\end{equation}
}
Applying the monotonicity of the soft Bellman operator repeatedly yields:
{\boldmath
\begin{equation}
Q^{\pi_{\text{new}}}(s, a) = \lim_{k \to \infty} (\mathcal{T}^{\pi_{\text{new}}})^k Q^{\pi_{\text{old}}}(s, a) \ge Q^{\pi_{\text{old}}}(s, a)
\end{equation}
}
This completes the proof of monotonic soft policy improvement.
\end{proof}

\begin{theorem}[Contraction Mapping of the Soft Bellman Operator]
The soft Bellman operator $\mathcal{T}^\pi$ defined on the Banach space $(\mathcal{B}(\mathcal{S} \times \mathcal{A}), \|\cdot\|_\infty)$ is a $\gamma$-contraction mapping.
\end{theorem}

\begin{proof}
Let $Q_1, Q_2 \in \mathcal{B}(\mathcal{S} \times \mathcal{A})$. For any state-action pair $(s, a)$:
{\boldmath
\begin{align}
|(\mathcal{T}^\pi Q_1)(s, a) - (\mathcal{T}^\pi Q_2)(s, a)| &= \gamma \left| \mathbb{E}_{s' \sim p} \left[ V_1(s') - V_2(s') \right] \right| \\
&\le \gamma \max_{s'} |V_1(s') - V_2(s')|
\end{align}
}
Evaluating the difference in soft value functions:
{\boldmath
\begin{align}
|V_1(s') - V_2(s')| &= \left| \mathbb{E}_{a' \sim \pi} \left[ Q_1(s', a') - \alpha \ln \pi(a'|s') - (Q_2(s', a') - \alpha \ln \pi(a'|s')) \right] \right| \\
&= \left| \mathbb{E}_{a' \sim \pi} \left[ Q_1(s', a') - Q_2(s', a') \right] \right| \\
&\le \mathbb{E}_{a' \sim \pi} \left[ |Q_1(s', a') - Q_2(s', a')| \right] \le \|Q_1 - Q_2\|_\infty
\end{align}
}
Taking the supremum over all $(s, a)$ yields:
{\boldmath
\begin{equation}
\|\mathcal{T}^\pi Q_1 - \mathcal{T}^\pi Q_2\|_\infty \le \gamma \|Q_1 - Q_2\|_\infty
\end{equation}
}
Because $\gamma \in (0, 1)$, $\mathcal{T}^\pi$ is a contraction mapping. By the Banach fixed-point theorem, repeated iteration converges to a unique soft fixed-point $Q^*$.
\end{proof}

\section{Foundational Architectural Questions and Epistemic Meta-Questions}

\subsection{Architectural Question 1: Fixed vs Automatic Temperature}
\textbf{Question:} Why is automatic temperature adjustment preferred over a fixed entropy constant $\alpha$?\\
\textbf{Answer:} Different environments feature vastly different reward scales and state complexities. A fixed $\alpha$ that works for a locomotion task may cause excessive randomness on a manipulation task or insufficient exploration on a navigation task. By formulating temperature tuning as a dual optimization problem against a target entropy $\bar{\mathcal{H}} = -\dim(\mathcal{A})$, SAC dynamically maintains the target entropy across the entire training lifecycle without manual tuning.

\textbf{Meta-Question:} How does maximum entropy reinforcement learning prevent policy collapse?\\
\textbf{Answer:} In standard deterministic policy gradients, when an action accidentally yields high reward, the policy rapidly concentrates probability mass into that single action, collapsing exploration. In max-ent RL, actions with equal or near-equal values receive spread-out probability mass according to the Boltzmann distribution, preserving exploration and resilience to environmental perturbations.

\subsection{Architectural Question 2: Clipped Double Q Target Strategy}
\textbf{Question:} Why does SAC use the minimum of two critic estimates for the target?\\
\textbf{Answer:} Function approximation introduces noise into value estimates. Because policy optimization inherently seeks high-value regions, errors compound positively, leading to runaway Q-value overestimation. Taking $\min(Q_1, Q_2)$ acts as a conservative bound, suppressing overestimation and stabilizing the off-policy learning loop.

\textbf{Meta-Question:} What is the mathematical connection between SAC and Energy-Based Models (EBMs)?\\
\textbf{Answer:} At the optimal soft policy fixed point, the policy takes the exact form of an energy-based Gibbs distribution $\pi^*(a | s) \propto \exp\left(\frac{1}{\alpha} Q^*(s, a)\right)$. Minimizing the KL divergence in policy optimization corresponds to minimizing the variational free energy of the policy distribution with respect to the energy landscape defined by the Q-critic.

\section{Algorithmic Implementation Specification}

\begin{algorithm}[H]
\caption{Soft Actor-Critic Evaluation and Dual Temperature Update}
\begin{algorithmic}[1]
\Require Transitions batch $(\mathbf{q}_1, \mathbf{q}_2, \mathbf{q}_{1,\text{targ}}, \mathbf{q}_{2,\text{targ}}, \mathbf{r}, \mathbf{lp}_{\text{next}}, \mathbf{d}, \mathbf{lp}_{\pi}, \mathbf{q}_{1,\pi}, \mathbf{q}_{2,\pi})$, $\alpha, \gamma, \bar{\mathcal{H}}, \eta_\alpha$
\Ensure Targets $\mathbf{y}$, critic loss $\mathcal{L}_Q$, actor loss $\mathcal{L}_\pi$, dual loss $\mathcal{L}_\alpha$, updated $\alpha'$
\State Initialize $\mathbf{y} \in \mathbb{R}^N, \quad err_1 \gets 0.0, \quad err_2 \gets 0.0$
\For{$i = 0$ \textbf{to} $N-1$}
    \If{$\mathbf{d}[i]$ is True}
        \State $\mathbf{y}[i] \gets \mathbf{r}[i]$
    \Else
        \State $min\_next \gets \min(\mathbf{q}_{1,\text{targ}}[i], \mathbf{q}_{2,\text{targ}}[i])$
        \State $\mathbf{y}[i] \gets \mathbf{r}[i] + \gamma \cdot (min\_next - \alpha \cdot \mathbf{lp}_{\text{next}}[i])$
    \EndIf
    \State $err_1 \gets err_1 + (\mathbf{q}_1[i] - \mathbf{y}[i])^2$
    \State $err_2 \gets err_2 + (\mathbf{q}_2[i] - \mathbf{y}[i])^2$
\EndFor
\State $\mathcal{L}_Q \gets (err_1 + err_2) / N$
\State $\mathcal{L}_\pi \gets \frac{1}{N}\sum_{i=0}^{N-1} (\alpha \cdot \mathbf{lp}_\pi[i] - \min(\mathbf{q}_{1,\pi}[i], \mathbf{q}_{2,\pi}[i]))$
\State $\bar{H} \gets - \frac{1}{N}\sum_{i=0}^{N-1} \mathbf{lp}_\pi[i]$
\State $\mathcal{L}_\alpha \gets - \frac{1}{N}\sum_{i=0}^{N-1} \alpha (\mathbf{lp}_\pi[i] + \bar{\mathcal{H}})$
\State $g_\alpha \gets - (-\bar{H} + \bar{\mathcal{H}})$
\State $\alpha' \gets \max(10^{-6}, \alpha - \eta_\alpha \cdot g_\alpha)$
\State \Return $(\mathbf{y}, \mathcal{L}_Q, \mathcal{L}_\pi, \mathcal{L}_\alpha, \alpha', \bar{H})$
\end{algorithmic}
\end{algorithm}

\section{Big-$\Theta$ Complexity Analysis}

\begin{itemize}
    \item \textbf{Soft Target Construction Time}: $\Theta(N)$ scalar operations across the batch.
    \item \textbf{Clipped Critic Error Time}: $\Theta(N)$ scalar multiplications and additions.
    \item \textbf{Policy and Dual Gradient Time}: $\Theta(N)$ evaluations of log probabilities.
    \item \textbf{Space Complexity}: $\Theta(N)$ memory buffer for temporal difference targets.
    \item \textbf{Parallel Reduction Depth}: $\mathcal{D} = \Theta(\log N)$ for batch mean computations.
\end{itemize}

\section{Single Canonical Repository Reference}

The complete deterministic scalar implementation, verification suite, and unit test benchmarks are published at:
\begin{center}
\url{https://github.com/Chief-Strategist-J/llm-obs-infra}
\end{center}

\end{document}
